% 09-evolution.tex — OpenBench Software Design Document. Written from .scratch/design-doc/brief.md on 2026-10-05.
\section{Maintenance and Evolution}\label{sec:evo}

Release~1 is fixed. Stretch items are built only if time allows, and Future items have no requirements yet. For each of them, and for the operational changes a successful deployment would bring, the design names the one place it lands, so none of them is a precondition for Release~1 and none changes the dependency direction of Section~\ref{sec:arch}.

\subsection{Extension Points}

Table~\ref{tab:evo-extensions} covers every Stretch and Future item of the SRS scope table plus two operational changes. Most rows add a class or a column behind an interface that already exists.

{\footnotesize\setlength{\tabcolsep}{4pt}\renewcommand{\arraystretch}{1.08}
\begin{longtable}{@{}L{0.25\textwidth}L{0.09\textwidth}L{0.115\textwidth}L{0.3\textwidth}L{\dimexpr0.245\textwidth-8\tabcolsep\relax}@{}}
\caption{Extension points for the Stretch, Future, and operational items.}\label{tab:evo-extensions}\\
\toprule
\textbf{Item} & \textbf{Tier} & \textbf{Where it lands} & \textbf{What changes} & \textbf{What does not change} \\
\midrule
\endfirsthead
\toprule
\textbf{Item} & \textbf{Tier} & \textbf{Where it lands} & \textbf{What changes} & \textbf{What does not change} \\
\midrule
\endhead
\midrule\multicolumn{5}{@{}r@{}}{\itshape continued on the next page}\\
\endfoot
\bottomrule
\endlastfoot
Service tiers and preemption on Metered pools (\req{Metered.Tier.Define}, \req{Metered.Tier.Bump}) & Stretch & \module{Allocation}, \module{Ledger}, \module{Settings}, \module{Requests} (a tier column and setter) & \code{MeteredAllocator} also implements \code{Bumpable}, planning Bumps of lower-tier units outside the bump cutoff; a tier premium multiplies the Unit Price in \code{Pricing.quote} as the priority premium does; tiers per type become settings; \code{request\_line} gains a tier & the lifecycle in \code{Request\-Lifecycle\-Service}, which already asks \code{AllocatorRegistry} for \code{Bumpable} on any class; the Bump Refund and Waitlist path; \code{ExclusiveAllocator} \\
Demand surcharge (\req{Ledger.Surcharge}) & Stretch & \module{Requests}, \module{Ledger}, \module{Settings} & \module{Requests} reads the requested window's fullness (allocated plus held over capacity) from \code{AllocationQueries} and passes it into \code{Pricing.quote} as an extra input; one multiplier inside \code{quote}; the surcharge per type is a setting & \code{Pricing} stays pure and \module{Ledger} never calls \module{Allocation}; the Quote still locks at submission and the Charge equals it \\
Actual usage recording and comparison (\req{Metered.Usage.Record}, \req{Metered.Usage.Compare}, \req{Dashboard.Usage}) & Stretch & \module{Allocation}, \module{Dashboard} & a \code{used\_units} column on \code{metered\_allocation}, written at Period end by a new Resource Owner operation in \module{Allocation}; one \module{Dashboard} view comparing requested and used units per Team and type & \module{Ledger}, since usage has no Ledger effect; the hourly check, which keeps using requested units \\
Email notifications (\req{CI-inapp}) & Stretch & \module{Notifications}, \module{Settings} & \code{EmailDelivery} implements \code{Delivery} and sends after commit; SMTP host, port, and sender become settings the Administrator edits, the SMTP password an environment variable & every caller of \code{Notifier}; the in-app row stays the record \\
Demand analytics on declined Suggestions & Future & \module{Dashboard} & one view over Declined \code{suggestion} rows (type, quantity, window, skills, reason) per Period & no data changes: Suggestions keep their original ask and are never deleted (\req{Suggest.Create}) \\
Person consent (lifts \req{LI-consent}) & Future & \module{Allocation} & \code{MatchedAllocator.select} first creates a consent approval addressed to the Person, then the Manager's; \code{approval} gains a kind; a Person's refusal takes the refusal path & \code{decideApproval}, the Objection Window, escalation, \module{Requests}, \module{Ledger} \\
Skill-based Day Rate & Future & \module{Settings}, \module{Ledger} & a rate per skill in Settings with the one Day Rate as the default; a rate lookup next to \code{Pricing} computes a Loan's Earning from its skills & \code{postEarning}, \code{correctEarning}, Loans, the Shortlist; Matched lines still cost the borrower nothing (\adr{0004}) \\
Team rank priority & Future & \module{Allocation}, \module{Settings} & a rank per Team set by the Capacity Planner; one comparator in the Bump plan of \code{ExclusiveAllocator} consults it & Priority by premium (\adr{0003}); the exclusion constraint; Refunds; \module{Requests} \\
Inter-Team transfers & Future & \module{Ledger} & a fifth Transaction kind, Transfer, and one operation posting a pair of Transactions between two Teams with a reason; both Statements show it; the glossary's ``exactly four kinds'' becomes five & append-only Transactions, Account arithmetic, the other four kinds, every caller of \code{LedgerService} \\
Forecasting accuracy reward & Future & \module{Ledger}, \module{Jobs} & at rollover, a Grant computed from requested versus used units (needs usage recording) & the Grant posting path and its Transaction kind \\
Usage feeds & Future & \module{Allocation} & an adapter fills \code{used\_units} from a scheduler or license server instead of the Resource Owner & as for usage recording \\
Auctions & Future & \module{Allocation} & a new Allocation Class whose allocator defers, as \code{ApprovalDriven} does, until bidding closes, with the winning bid as the Quote; the largest Future item, needing its own design pass & the four Transaction kinds (a winning bid posts a Charge); the existing allocators \\
SSO (lifts \req{LI-integrations}) & Future & \code{security} & a Spring Security authentication provider (OpenID Connect or SAML) mapping the identity provider's subject to \code{member.username}; password login can be switched off & \module{Members}, roles, sessions, every \code{@PreAuthorize} \\
Calendar integration (lifts \req{AS-calendar}) & Future & \module{Members} & an adapter behind \code{isAvailable} and \code{availableWithSkills} that also removes calendar busy days & \module{Allocation} and the Shortlist call the same operations \\
HR integration & Future & \module{Members} & an import adapter that creates, moves, and deactivates Members and Teams from the HR system & every module above \module{Members}; the Administrator screens stay for corrections \\
Cluster-scheduler integration & Future & \module{Inventory}, \module{Jobs} & an adapter behind \module{Inventory} sets Resource units from the scheduler through \code{setUnits} and its guard; a new consumer of allocation records, a polling step in \module{Jobs}, pushes committed Metered allocations to the scheduler & allocation transactions make no remote call; \code{MeteredAllocator}; there is still no event bus \\
A second application instance & Oper\-ational & \module{Jobs}, \code{security} & a job lock table each job claims with a conditional update before it runs; sticky sessions at nginx or a shared session store (Spring Session over JDBC) & allocation correctness: type locks and exclusion constraints live in the database and already hold across instances \\
A new Allocation Class (\req{QA-extensible}, \req{CO-plugin}) & Oper\-ational & \module{Allocation} & one \code{Allocator} bean and the capability interfaces it needs; its name in the \code{AllocationClass} enumeration and the check on \code{resource\_type} & \module{Requests}, \module{Ledger}, and the existing allocators: \code{AllocatorRegistry} finds the bean by \code{allocationClass()} \\
\end{longtable}
}

\textbf{One dependency added during design.} \module{Ledger} calls \module{Notifications} to tell each Manager about Period rollover and mid-Period Grants (Table~\ref{tab:cross-notify}). \module{Notifications} sits in the bottom layer, so the dependency rule still holds and the ArchUnit test allows it.

\subsection{Assumptions That Would Break the Design}

Table~\ref{tab:evo-assumptions} lists the assumptions the design leans on hardest, what relies on each, and what would have to change if it stopped being true.

{\footnotesize\setlength{\tabcolsep}{4pt}\renewcommand{\arraystretch}{1.08}
\begin{longtable}{@{}L{0.22\textwidth}L{0.37\textwidth}L{\dimexpr0.41\textwidth-4\tabcolsep\relax}@{}}
\caption{Assumptions and what depends on them.}\label{tab:evo-assumptions}\\
\toprule
\textbf{Assumption} & \textbf{What relies on it} & \textbf{If it changed} \\
\midrule
\endfirsthead
\toprule
\textbf{Assumption} & \textbf{What relies on it} & \textbf{If it changed} \\
\midrule
\endhead
\bottomrule
\endlastfoot
One company in one time zone (\req{LI-company}) & Price Bands, Period boundaries, person-days, the rollover time, and the API's offsets all use one configured zone & instants are already stored zone-free, but Price Band lookup, Periods, and rollover would need a zone per Team or Resource, and Periods would no longer turn over at one moment \\
One server instance, no high availability (\req{LI-company}, \req{QA-available}) & each job runs once because one process runs it; sessions live in memory & a job lock table and shared sessions (Table~\ref{tab:evo-extensions}); allocation correctness is unaffected \\
Hour granularity for Metered pricing and checks; day granularity for Loans & the hourly sum and the Quote step through whole hours; Loans, Earnings, and the \code{loan} exclusion constraint use dates & finer windows multiply the hourly check's buckets and change Quote arithmetic; half-day Loans change person-days, the Day Rate, and the date-range constraint \\
The Capacity Planner sets Grants before the first Period (\req{AS-planner}) & every Exclusive and Metered line needs credits; rollover posts the planned Grants & without Grants only Matched lines can be submitted (\code{INSUFFICIENT\_BALANCE}); the design invents no default Grant \\
Loans are recorded honestly (\req{AS-honest}) & the Earning posts when the Loan forms and is corrected only by an early end & confirming the work would need a record of days actually worked, as usage recording does for Metered pools, and an Earning posted on confirmation \\
The scale of \req{PE-scale}: 500 Members, 200 Resources, 10,000 Requests a year & the Shortlist is built on read, hourly sums are computed by query, and the Waitlist is walked inside the releasing transaction & at a hundred times the scale: stored hourly totals per type, a precomputed Shortlist, and a Waitlist fill moved to \code{DueWorkJob} \\
\end{longtable}
}

\subsection{Risks}

Table~\ref{tab:evo-risks} lists the technical and schedule risks of this design; likelihood and impact are low, medium, or high, as in the SRS.

{\footnotesize\setlength{\tabcolsep}{4pt}\renewcommand{\arraystretch}{1.08}
\begin{longtable}{@{}L{0.2\textwidth}L{0.07\textwidth}L{0.15\textwidth}L{\dimexpr0.58\textwidth-6\tabcolsep\relax}@{}}
\caption{Risks and their mitigation in the design.}\label{tab:evo-risks}\\
\toprule
\textbf{Risk} & \textbf{Likeli\-hood} & \textbf{Impact} & \textbf{Mitigation in the design} \\
\midrule
\endfirsthead
\toprule
\textbf{Risk} & \textbf{Likeli\-hood} & \textbf{Impact} & \textbf{Mitigation in the design} \\
\midrule
\endhead
\bottomrule
\endlastfoot
Lock-order discipline breaks: a new path consumes capacity without the type locks, or takes them out of order & medium & high: an oversold Metered pool, or deadlocks & \code{AllocationLock.lockTypes} takes a sorted set, so the order holds by construction; every path that consumes or reduces capacity (submit, Waitlist fill, Bump, Pool Capacity reduction) runs through \code{RequestLifecycleService} or \code{ReductionGuard}, which lock first; the exclusion constraints backstop Exclusive instances and Loans; \code{lock\_timeout} turns a stuck wait into a 503. Metered pools have no constraint backstop, so the Metered race test and code review guard them \\
The exclusion constraints depend on the \code{btree\_gist} extension & low & high: the server does not start & the first Flyway migration creates the extension, so a missing privilege fails at start-up and not at the first Request; the extension ships with PostgreSQL's contributed modules and the official image; the installation guide states what the database role needs (\req{UD-admin}); every test run proves it on PostgreSQL~16 \\
Waitlist walk cost grows with the Waitlist & low at \req{PE-scale} & medium: Releases slow down and hold the type lock longer & \code{waitingFor(type, freed)} reads only entries whose window overlaps the freed one, through an index on type and start; expired entries leave through \code{expireWaitingLines}; the walk stops at the first Request that allocates; logged durations show growth; if needed, the fill moves to \code{DueWorkJob} and still meets \req{PE-jobs} \\
The single polling job stalls or fails & low & medium: lapses and expiries come late (\req{PE-jobs}) & one transaction per Request, so one failure does not stop the rest; due work is found by query, so a missed run or a restart is caught up on the next run; a scheduler pool of two threads, so rollover never delays \code{DueWorkJob}; one log line per run with counts and duration \\
The front end drifts from the API & medium & medium: screens fail at run time & client types are generated from the OpenAPI document on every build, so a changed endpoint fails the TypeScript compile; the end-to-end test; Quotes and rules live only on the server \\
Four students, one semester (\req{RI-5}) & high & high: Release~1 incomplete & Release~1 is fixed and every Stretch and Future item is an extension point, never a precondition; one process, one database, no broker; modules are built bottom-up in the layer order against the interfaces of Section~\ref{sec:modules}, each testable alone; Stretch work starts only after the end-to-end test of the three-line Request passes \\
\end{longtable}
}
